%% GENERATED by python/paper/build.py from results/ -- do not edit by hand.
%% Rebuild:  .venv\Scripts\python.exe python\paper\build.py --only OECD_Sources
%% Source:   data/oecdFigure1_sources.csv, data/oecd26_theta.csv, data/oecd26_pensionTiers_harmonised.csv, data/oecd29_sources.csv
\begin{table}[!htb]
\centering
\footnotesize
\begin{threeparttable}
\caption{The sources of figure \ref{fig:US:OECDdata}}
\label{table:OA:OECD_Sources}
\renewcommand{\arraystretch}{1.25}
\begin{tabularx}{\textwidth}{>{\raggedright\arraybackslash}p{3.8cm}>{\raggedright\arraybackslash}X>{\raggedright\arraybackslash}X>{\raggedright\arraybackslash}X}
\toprule
\textbf{Columns} & \textbf{Source} & \textbf{Series and years} & \textbf{Note} \\
\midrule 
\multicolumn{4}{l}{\textit{Plotted in the figure of the paper}} \\
$\theta$, the design, and its index & OECD, Pensions at a Glance 2021, Table 4.2 Gross pension replacement rates from mandatory public, mandatory private and voluntary private schemes (StatLink uwc8nt, Version 1 - Last updated: 08-Dec-2021); retrieved 2026-10-30. Top-ups and tested benefits: OECD SOCX 2020 and national sources (data/oecd26\_topups.csv; python/paper/oecd26Theta.py, S\_TOPUP) & $\theta$ is the earnings-linked share of the least-squares line through the public pension at 0.5, 1 and 2 times average earnings ($c/(a+c)$ of $p(y) = a + c\,y$), times one minus the share of public spending on minimum top-ups and tested benefits: the gross public pension of men with a full career under the rules of 2020 (lowest earnings: New Zealand at 0.6), public spending as pensionSpending. The index is $1/(2-\theta)$. Years: 2020 (21 of 21) & The paper\textquotesingle s measure (appendix \ref{app:oecd26:data}), variant C of data/oecd26\_theta.csv, on the 21 countries where funded tiers pay less than a quarter of mandatory benefits in 2020 and the public pension is not tested against them. The UK\textquotesingle s $\theta$ is not read off the benefit formula: it is the State Second Pension\textquotesingle s share of spending on the State Pensions and Pension Credit in 2019--20. \\ % row: theta
pensionSpending & OECD: Social Expenditure Database, OECD.ELS.SPD:DSD\_SOCX\_AGG@DF\_SOCX\_AGG(1.0); Data Explorer; retrieved 2026-09-30 & public (ES10), old age and survivors (TP01), in cash (C), \% of GDP. Years: 2020 (29 of 29) & The OECD\textquotesingle{}s \textquotesingle{}pension spending\textquotesingle{}: cash expenditure on old-age and survivors pensions (Pensions at a Glance measure PEP to its rounding). \\ % row: pensionspending
pop1990\_oecd, pop2020\_oecd, popGrowth & OECD: Historical population data, OECD.ELS.SAE:DSD\_POPULATION@DF\_POP\_HIST(1.0); Data Explorer; retrieved 2026-09-30 & POP, persons, both sexes, all ages, historical; popGrowth = pop2020/pop1990. Years: 1990 and 2020 (29 of 29) & France: the series adds the overseas departments from 1991, so its ratio (1.192) is not on a constant territory (World Bank 1.160). \\ % row: pop1990_oecd
yHighOverLow, yHighOverLow\_year & OECD: Earnings database, Decile ratios of gross earnings, OECD.ELS.SAE:DSD\_EARNINGS@DEC\_I(1.0); Data Explorer; retrieved 2026-10-23 & y3/y1 of the three income groups at the U.S. population shares 0.582/0.176/0.242 (data/oecd29\_distribution.csv, role data, source oecd\_earnings\_men\_ft, convention usPct); men (SEX=M), full-time employees, D5\_1 = P50/P10 and D9\_5 = P90/P50; two-piece lognormal through them. Years: 2019 (20 of 21); 2018: POL & Not a column of data/oecdFigure1.csv: derived at build time from the income groups the model is fed (python/US/oecd26.py), sourced in data/oecd29\_sources.csv. \\ % row: yhighoverlow
\multicolumn{4}{l}{\textit{Further concepts in data/oecdFigure1.csv}} \\
pensionSpending\_cashInKind & OECD: Social Expenditure Database, OECD.ELS.SPD:DSD\_SOCX\_AGG@DF\_SOCX\_AGG(1.0); Data Explorer; retrieved 2026-09-30 & as pensionSpending, cash plus in kind (\_T). Years: 2020 (29 of 29) & The in-kind part is services to the elderly (care, home help), largest in NOR 2.57, SWE 2.43, DNK 1.87, FIN 1.61. \\ % row: pensionspending_cashinkind
pop1990\_wb, pop2020\_wb, popGrowth\_wb & World Bank: World Development Indicators; API lastupdated 2026-07-13; retrieved 2026-09-30 & SP.POP.TOTL; popGrowth\_wb = pop2020/pop1990. Years: 1990 and 2020 (29 of 29) & -- \\ % row: pop1990_wb
grr050\_men, grr100\_men, grr050\_women, grr100\_women & OECD: Pensions at a Glance 2021, table 4.1; StatLink xlsx ``Version 1 - Last updated: 08-Dec-2021''; retrieved 2026-09-30 & gross pension replacement rate at 0.5 and 1.0 times average earnings, mandatory schemes, \% of individual earnings, one decimal as published; women = men where the table gives none. Years: pension rules in 2020 (29 of 29) & The OECD Data Explorer dataflows DSD\_PAG@DF\_PRR and DSD\_PAG@DF\_PAG now hold only 2024 (Pensions at a Glance 2025); the publication StatLink still serves the 2020 vintage. \\ % row: grr050_men
bbIndex & derived: Pensions at a Glance 2021, table 4.1 & grr100\_men / grr050\_men. Years: 2020 (29 of 29) & Men: the replacement-rate ratio of the table headline, not the paper\textquotesingle{}s Bismarckian index (data/oecd26\_theta.csv). \\ % row: bbindex
bbIndex\_women & derived: Pensions at a Glance 2021, table 4.1 & grr100\_women / grr050\_women. Years: 2020 (29 of 29) & -- \\ % row: bbindex_women
bbIndex\_mean & derived: Pensions at a Glance 2021, table 4.1 & (bbIndex + bbIndex\_women) / 2. Years: 2020 (29 of 29) & -- \\ % row: bbindex_mean
gini, gini\_year, gini\_wiidId, gini\_source & UNU-WIDER: World Income Inequality Database (WIID), version 8 September 2026, doi 10.35188/UNU-WIDER/WIID-080926; version 8 Sep 2026; retrieved 2026-09-30 & the observation WIID flags for its Companion (wiidcompanion = Yes): disposable income, persons, per capita except where gini\_source says otherwise; 0-1 scale. Years: 2020 (29 of 29) & gini\_wiidId is the WIID row id. \\ % row: gini
gini\_wiidMarket, gini\_wiidMarket\_year & UNU-WIDER: WIID, version 8 September 2026; version 8 Sep 2026; retrieved 2026-09-30 & source OECD.Stat (Income Distribution Database), market income, square-root scale, persons; 0-1 scale. Years: 2020 (22 of 29); 2018: JPN; 2019: DNK, FRA, DEU, CHE, TUR; none in 2018-2022: ISL & Before taxes and transfers; retirees without market income count at zero. \\ % row: gini_wiidmarket
gini\_widPretax, gini\_widPretax\_year & World Inequality Lab: World Inequality Database (WID), bulk download; Last-Modified Wed, 09 Sep 2026 14:06:24 GMT; retrieved 2026-09-30 & gptincj992, p0p100: Gini of pre-tax national income, equal-split adults (20+). Years: 2020 (29 of 29) & -- \\ % row: gini_widpretax
\bottomrule
\end{tabularx}
\begin{tablenotes}[flushleft]
\footnotesize
\item[] \textit{Note:} The columns are those of data/oecdFigure1.csv, but for $\theta$ and its index, from data/oecd26\_theta.csv (data/README.md), and yHighOverLow, derived from the income groups of data/oecd29\_distribution.csv; data/oecdFigure1\_sources.csv, data/oecd26\_pensionTiers\_harmonised.csv, data/oecd29\_sources.csv and the files the $\theta$ row names carry the URL of every source. bbIndex, the ratio of the replacement rates at average and at half of average earnings, is not the paper\textquotesingle s index.
\end{tablenotes}
\end{threeparttable}
\end{table}
